\subsection{An entire class-two two-primary tower and every odd pro-prime kernel}

\begin{statement}{D.29}{Proved}{A compatible infinite nonabelian lift of the alternate CM pair}
For $k\geq1$ define the characteristic quotients
\[
 H^{(k)}=B/\bigl(\gamma_3N\,[N,N]^{2^k}N^{2^{k+1}}\bigr),
 \qquad |H^{(k)}|=576\cdot2^{4753k+97}.                  \tag{D.29.1}
\]
Thus $H^{(1)}=H_2$ of D.27. The alternate CM pair
$A=I$, $\beta=\left(\begin{smallmatrix}9&8\\4&9\end{smallmatrix}\right)$
has actual cusp-preserving automorphisms on $H^{(2)}$
and $H^{(3)}$. Their corrected cusp words have lengths
1816936 and 4448112, respectively, and are checked in the
complete truncated Magnus models over $\mathbb Z/4$
and $\mathbb Z/8$. Their multipliers are 1 and 289 modulo
288 and 576. They retain D.27's homology data, including
its right $s$-centralizer correction, while changing
the central coordinates as explicitly recorded.

There is a compatible infinite lift to
\[
 \mathcal P_2^{[2]}:=\varprojlim_k H^{(k)},\qquad
 1\longrightarrow N^{(2)}/\overline{\gamma_3N^{(2)}}
 \longrightarrow\mathcal P_2^{[2]}
 \longrightarrow Q\longrightarrow1.                  \tag{D.29.2}
\]
It induces the particular $H^{(1)}$ action obtained by
reducing the recorded $\mathbb Z/4$ witness. This is a
specified replacement of D.27's central representative
within the same homology class; it is not asserted to
extend every previously recorded central representative.

Let $\mathcal P_{\rm odd}$ be the fibre product over $Q$
of the relative completions $\mathcal P_p$ for every odd
prime $p$. The new action extends to an automorphism of
\[
 \mathcal R=\mathcal P_2^{[2]}\times_Q\mathcal P_{\rm odd},
 \qquad
 \ker(\mathcal R\to Q)=
 \bigl(N^{(2)}/\overline{\gamma_3N^{(2)}}\bigr)
 \times\prod_{p\text{ odd}}N^{(p)}.                    \tag{D.29.3}
\]
Its action on every odd relative completion is the identity,
with the prescribed conjugators at three. Its multiplier
has $\lambda_2\in1+16\mathbb Z_2$ and $\lambda_p=1$
for every odd $p$. A single profinite cusp gauge supplies
compatible nonarithmetical shadows on every finite
characteristic refinement retaining its $G_6$ and CM data.

The mixed-prime witness $J$ of D.28 agrees with this action
on the exact common quotient $C=H\times_QG_3$ constructed
there. Hence it further extends to
\[
                J\times_C\mathcal R.                 \tag{D.29.4}
\]
The realization of this CM pair by an automorphism of the
full $\widehat B$, and hence by genuine profinite GT$_0$,
remains undecided. No replacement of the all-finite-quotient
test of D.22 by (D.29.2)--(D.29.4) is made.
\end{statement}
\begin{proof}
The truncated Magnus model over $\mathbb Z/2^k$ has
multiplication
\[
 (v,w)(v',w')=(v+v',w+w'+v(v')^{\mathsf T}).             \tag{D.29.5}
\]
The image of a free generator is $(e_i,0)$, its inverse
is $(-e_i,e_ie_i^{\mathsf T})$, and its $a$th power is
$(ae_i,\binom a2 e_ie_i^{\mathsf T})$. The pair of
diagonal coordinates $a\bmod2^k$ and $\binom a2\bmod2^k$
determines $a\bmod2^{k+1}$: replacing $a$ by $a+2^k$
changes the second coordinate by $2^{k-1}$ modulo $2^k$.
The independent off-diagonal commutators have order $2^k$.
Collection in a free class-two group therefore proves
faithfulness, the kernel in (D.29.1), and its order
$2^{(k+1)97+k\binom{97}{2}}$ over $Q$.
The graph-word evaluation of D.27 applies with integer
coefficients modulo $2^k$, so it computes actual words
in these quotients.

First consider the lift from $\mathbb Z/2$ to
$\mathbb Z/4$. Inside the latter model, let $C_1$ have
basis $a_i^2$ and $[a_i,a_j]$. It is an abelian group
$(\mathbb Z/4)^{4753}$: the commutator of two squares
is a fourth power of a basic commutator, which vanishes.
Its elements are identified by the upper-triangular
Magnus coefficients, with $v_i=2w_{ii}$ and
$w_{ji}=-w_{ij}$ for $i\ne j$. Its subgroup $2C_1$ is
the new kernel over $H^{(1)}$.

Fix the full D.27 conjugators, including their recorded
central and right-centralizer words, and write
$g_s=\beta s\beta^{-1}$, $g_r=ArA^{-1}$.
Conjugation gives integral linear automorphisms $S,R$ of
$C_1$. Changing the conjugators on the left by $b,a\in C_1$
changes the cusp error by
\[
       (1-S)b+S(1-R)a-hc_1,                            \tag{D.29.6}
\]
where $c_1$ is the $C_1$ coordinate of $t^{72}$ and the
multiplier is changed by $72h$. This formula is exact
because $C_1$ is abelian; it includes the actual $g_s,g_r$
actions rather than replacing them by their $Q$ images.
Modulo two its matrix $L_0$ has rank 4738, as in D.27,
and its kernel has dimension $9507-4738=4769$.

The code forms all fifteen exact dual functionals on
$\operatorname{coker}L_0$. For a relation $z\in\ker L_0$
it computes the new obstruction
\[
 z\longmapsto\tfrac12 Lz\pmod{\operatorname{im}L_0},
 \qquad
 \ker L_0\longrightarrow\operatorname{coker}L_0.        \tag{D.29.7}
\]
Choosing a different zero-one lift of $z$ changes this
class by an $L_0$ image, so the map is well-defined.
The exact rank is 15. Its saved relation certificates,
followed by an ordinary $L_0$ correction, solve the full
$\mathbb Z/4$ equation. Generation is unchanged on
$N^{\rm ab}/2$, where the rank remains97; the finite
two-kernel Frattini argument gives an automorphism.
Direct evaluation of the resulting 1816936-letter cusp
word has zero in every degree-one and degree-two
coordinate. These are the actual results of
\texttt{d\_cm\_exponent4.py}, with all coefficients saved
in its JSON output.

We prove the infinite extension by exact stabilization.
For $k\geq1$ put
\[
 C_k=\ker(H^{(k+1)}\to H^{(k-1)}),
 \quad H^{(0)}=G_2,
 \quad
 C_k=\langle a_i^{2^k},[a_i,a_j]^{2^{k-1}}\rangle
       \simeq(\mathbb Z/4)^{4753}.                     \tag{D.29.8}
\]
The commutator of two displayed powers has exponent
$2^{2k}$ and vanishes modulo $2^{k+1}$, so $C_k$ is
abelian. Its subgroup $2C_k$ is exactly
$\ker(H^{(k+1)}\to H^{(k)})$. Its conjugation action
factors through $G_2$: commutators of an element of
$N^2[N,N]$ with its generators vanish in this quotient.
Thus $C_k$ is here a $G_2$-module; the reduced module
$C_k/2C_k$ is the $Q$-module used in D.27.

Here are the formulas proving stabilization, including
the normalization of its diagonal coordinates. Suppose
the image of $a_i$ under conjugation by a fixed word is
$(v_i,W_i)$ in the Magnus model. Its $2^k$th power has
degree-two coordinate
$2^kW_i+\binom{2^k}{2}v_iv_i^{\mathsf T}$.
After division by $2^{k-1}$ its coordinate is
\[
              2W_i+(2^k-1)v_iv_i^{\mathsf T}\pmod4.   \tag{D.29.9}
\]
The diagonal coordinate must be multiplied by the inverse
of $2^k-1$ modulo four to express it in the basis
$a_j^{2^k}$. The image of a basic commutator is
$v_iv_j^{\mathsf T}-v_jv_i^{\mathsf T}$.
For $k\geq2$ the factor $2^k-1$ is always3 modulo four.
The vectors $v_i\bmod4$ depend only on the fixed image
in $Q$, and $W_i\bmod2$ depends only on the fixed
conjugator image in $G_2$. Consequently these matrices
are identical for every $k\geq2$.

The multiplier column stabilizes as well. Changing
$\lambda$ by $72\cdot2^{k-1}h$ changes the cusp power
by $(t^{12})^{3\cdot2^k h}$. If $t^{12}$ has Magnus
coordinates $(v,W)$, its normalized coordinate for
$k\geq2$ is
\[
                  2W+vv^{\mathsf T}\pmod4,             \tag{D.29.10}
\]
again with diagonal multiplier3. This follows by the
power formula in (D.29.5). In particular the multiplier
retains $\lambda\equiv1\pmod{72}$ at every finite stage.

The program \texttt{d\_cm\_pro2\_class2.py} computes this
one constant matrix modulo four and its map (D.29.7).
The latter again has rank 15, with explicit basis
certificates. Hence its image contains twice the entire
target module. To verify this last implication, let
$2e$ be any target. Surjectivity of (D.29.7) supplies
$z\in\ker L_0$ with $Lz/2-e\in\operatorname{im}L_0$;
an additional correction $2z'$ then makes
$L(z+2z')=2e$ modulo four.
The program also verifies an actual next lift: a
4448112-letter cusp word vanishes in the complete
$\mathbb Z/8$ Magnus model. This yields the stated
multiplier289 and kernel order $2^{14356}$.

We now construct a compatible infinite family, rather
than infer one from a few finite examples. Start with
the recorded $H^{(2)}$ action. Lift its two conjugators
and multiplier arbitrarily one level higher; its cusp
error lies in $2C_2$. The surjectivity just proved
corrects that error by elements of $C_2$ and the
specified multiplier increment. These corrections
fix the action on $H^{(1)}$. More generally, at stage
$k$ the same constant matrix corrects the error in
$2C_k$, preserves the action on $H^{(k-1)}$, and
produces an automorphism of $H^{(k+1)}$. The homology
parities never change, so the stabilization hypotheses
remain satisfied at every step. Although the chosen
action at level $k$ can change once, every level $k-1$
has already stabilized. Taking these stabilized
reductions gives compatible actions, conjugators and
multiplier on the entire inverse system. The multiplier
increments tend to zero two-adically. Every finite
action is an automorphism by the same rank 97 Frattini
argument, so their inverse limit is an automorphism.

The characteristic quotients in (D.29.1) are cofinal
among the finite two-group quotients of the free
class-two completion of $N$. This proves the exact
kernel and quotient map in (D.29.2). The construction
fixes the reduction of the recorded $H^{(2)}$ witness
to $H^{(1)}$; it does not claim that every central
coordinate of the $H^{(2)}$ witness remains fixed.
Both finite representatives are retained in the
output files, so the comparison is explicit.

For every odd $p$, the prescribed image $q(\beta)$
centralizes $s$ in $Q$. It lifts to
$C_{\mathcal P_p}(s)$ by coprime order-two conjugacy,
using the proof in D.27. At three retain the particular
$G_3$ conjugator there; at the other odd primes lift
its $Q$ image. Take $A=1$ and the identity action on
each $\mathcal P_p$. These actions have the same
elliptic reductions and cusp multiplier component1.
Finite kernels of distinct prime-power order have no
common nontrivial quotient, so the map from $\widehat B$
to their fibre product with (D.29.2) is onto. Passing
to all finite subproducts proves the exact quotient
map and kernel in (D.29.3).
The two-primary multiplier is 1 modulo 16; its odd
components are 1. This supplies a single
$\lambda\in\widehat{\mathbb Z}^{\times}$ agreeing on
all common finite cusp quotients.

Choose one lift of the combined $\beta$ to $\widehat F$.
The cusp gauge of D.26 then makes its associated $f$
a commutator and supplies both hexagons at every finite
level. It fixes the CM coordinates. They still disagree
on the common field of D.25, so every member retaining the $G_6$ and CM data
is nonarithmetical. D.28's additional mixed-prime
quotient was verified with these same homology parities,
including the right-centralizer correction, and the
same $G_3$ conjugators. Its exact common quotient with
(D.29.3) is the $C$ computed there. Thus the two actions
agree on $C$ and define (D.29.4). This last step uses
the specified quotient maps and compatibility, not
only the equality of their CM matrix pairs.
\end{proof}
